\chapter{The Actor Model \& Distributed Clustering (\texttt{std:actor} \& \texttt{std:actor:cluster})}

\section{What is an Actor?}

\begin{kidbox}{The Worker with the Private Office and Mailbox}
Imagine a smart worker named Sam who sits in a private office.
The office door is locked from the inside! Nobody is allowed to walk in and mess with Sam's desk. That means nobody can ever mess up Sam's toys or papers!
Outside the office door, there is a mail slot that drops into a locked \textbf{Mailbox}.
\begin{itemize}
    \item If you want Sam to do something, you write a letter and drop it into the mailbox.
    \item Sam opens letters one at a time, in the exact order they arrived (first in, first out).
    \item Sam can do some work, update notes on the desk, send a letter back to you, or hire a new worker!
\end{itemize}
This is the \textbf{Actor Model}! It makes writing concurrent software impossible to mess up with race conditions!
\end{kidbox}

\begin{syntaxbox}{Spawning Local Actors and Message Passing}
In Shae, you create an actor by providing an initial state and a message handler function:
\begin{lstlisting}[language=shaelang]
// 1. Define the initial state (bank balance)
let initial_balance = 100;

// 2. Spawn the actor
let bank_actor = actor.spawn(initial_balance, fn(balance, msg) {
    match msg {
        {"type": "deposit", "amount": amt} => {
            let new_bal = balance + amt;
            return Result.Ok(new_bal);
        },
        {"type": "withdraw", "amount": amt} => {
            if amt > balance {
                return Result.Err("Insufficient funds!");
            }
            return Result.Ok(balance - amt);
        },
        {"type": "get"} => {
            return Result.Ok(balance);
        },
        _ => return Result.Err("Unknown command"),
    }
});

// 3. Send asynchronous messages (tell)
bank_actor.tell({"type": "deposit", "amount": 50});

// 4. Send request-reply query (ask)
let current_funds = bank_actor.ask({"type": "get"});
println("Current Funds: " + current_funds); // 150
\end{lstlisting}
\end{syntaxbox}

\section{Distributed Clustering over Encrypted TLS Mesh}

\begin{kidbox}{The Global Pneumatic Post Office}
What if you have workers sitting in offices in London, New York, and Tokyo?
You can't just slide letters under the door!
Instead, they connect their offices with high-speed, underground, armored pneumatic tubes!
Every letter is locked with a secret cryptographic golden key so nobody on the street can read it!
In Shae, \texttt{std:actor:cluster} connects servers across the entire world into an encrypted supercomputer mesh!
\end{kidbox}

\begin{deepbox}{Cluster Architecture, TLS, and Wire Protocols}
Distributed Shae clusters (\texttt{src/actor\_cluster.rs}) establish a full-mesh topology across participating nodes:
\begin{itemize}
    \item \textbf{Transport Layer Security}: All node-to-node TCP links are encrypted using modern TLS (TLS 1.3 with AES-GCM or ChaCha20-Poly1305) via the \texttt{rustls} v0.23 library.
    \item \textbf{Mutual TLS (mTLS)}: Nodes require bidirectional cryptographic certificate authentication against a shared Root Certificate Authority (CA) bundle.
    \item \textbf{Wire Framing}: Messages are serialized into length-delimited binary frames with a 32-bit big-endian length prefix and a JSON/binary payload representing \texttt{ClusterWireMsg}:
    \begin{itemize}
        \item \texttt{ClusterWireMsg::Tell \{ target, msg \}}
        \item \texttt{ClusterWireMsg::Ask \{ target, msg, correlation\_id \}}
        \item \texttt{ClusterWireMsg::Reply \{ correlation\_id, result \}}
        \item \texttt{ClusterWireMsg::Reload \{ target, handler\_code, has\_migration \}}
    \end{itemize}
\end{itemize}
\end{deepbox}

\section{Hot-Code Reloading: Changing Tires at 200 MPH}

\begin{kidbox}{The Race Car Pit Stop Without Stopping}
Imagine you are driving a race car around a track at 200 miles per hour!
Suddenly, it starts raining, and you need wet-weather rain tires!
Do you stop the car, turn off the engine, tow the car back to the garage, and spend three weeks rebuilding it?
No! What if a superhero mechanic zoomed alongside your car and swapped all four tires while the wheels were still spinning at full speed?
In Shae, \textbf{Hot-Code Reloading} lets you upgrade the code inside running server actors without shutting them down, without losing a single unread message, and without breaking network connections!
\end{kidbox}

\begin{lstlisting}[language=shaelang]
// Upgrading the bank actor with an interest calculation feature
let new_handler = fn(balance, msg) {
    match msg {
        {"type": "add_interest", "rate": r} => {
            return Result.Ok(balance + (balance * r));
        },
        // Retain existing operations...
        _ => return Result.Ok(balance),
    }
};

// Upgrade the running actor instantly!
bank_actor.reload(new_handler);
\end{lstlisting}
